\section{Resultados: Polígonos PVA}

Las Figs.~\ref{fig:pva_0_1}--\ref{fig:pva_4_5} muestran el polígono de velocidades admisibles de cada robot en un instante representativo de la simulación, capturado cuando los robots se encuentran en la región de mayor proximidad entre sí.

Cada gráfica muestra:
\begin{itemize}
  \item \textbf{Región verde}: el polígono $\mathcal{P}(v,\omega)$ de velocidades admisibles, acotado por las restricciones del LiDAR y los límites de velocidad.
  \item \textbf{Punto azul} ($u_{\text{ref}}$): la velocidad de referencia generada por la ley de control nominal.
  \item \textbf{Punto rojo} ($u^*$): la velocidad óptima resultante de la proyección al polígono.
\end{itemize}

Cuando $u_{\text{ref}}$ cae dentro del polígono, se tiene $u^* = u_{\text{ref}}$ y no hay corrección por evasión. Cuando $u_{\text{ref}}$ cae fuera, $u^*$ es la proyección más cercana sobre la frontera del polígono, demostrando la acción correctiva del PVA.

\begin{figure}[H]
  \centering
  \begin{minipage}[b]{0.48\columnwidth}
    \centering
    \includegraphics[width=\linewidth]{pva_robot0}
    \subcaption{robot0}
  \end{minipage}
  \hfill
  \begin{minipage}[b]{0.48\columnwidth}
    \centering
    \includegraphics[width=\linewidth]{pva_robot1}
    \subcaption{robot1}
  \end{minipage}
  \caption{Polígonos PVA para robot0 y robot1.}
  \label{fig:pva_0_1}
\end{figure}

\begin{figure}[H]
  \centering
  \begin{minipage}[b]{0.48\columnwidth}
    \centering
    \includegraphics[width=\linewidth]{pva_robot2}
    \subcaption{robot2}
  \end{minipage}
  \hfill
  \begin{minipage}[b]{0.48\columnwidth}
    \centering
    \includegraphics[width=\linewidth]{pva_robot3}
    \subcaption{robot3}
  \end{minipage}
  \caption{Polígonos PVA para robot2 y robot3.}
  \label{fig:pva_2_3}
\end{figure}

\begin{figure}[H]
  \centering
  \begin{minipage}[b]{0.48\columnwidth}
    \centering
    \includegraphics[width=\linewidth]{pva_robot4}
    \subcaption{robot4}
  \end{minipage}
  \hfill
  \begin{minipage}[b]{0.48\columnwidth}
    \centering
    \includegraphics[width=\linewidth]{pva_robot5}
    \subcaption{robot5}
  \end{minipage}
  \caption{Polígonos PVA para robot4 y robot5.}
  \label{fig:pva_4_5}
\end{figure}

\begin{figure}[H]
  \centering
  \begin{minipage}[b]{0.48\columnwidth}
    \centering
    \includegraphics[width=\linewidth]{pva_robot6}
    \subcaption{robot6}
  \end{minipage}
  \hfill
  \begin{minipage}[b]{0.48\columnwidth}
    \centering
    \includegraphics[width=\linewidth]{pva_robot7}
    \subcaption{robot7}
  \end{minipage}
  \caption{Polígonos PVA para robot6 y robot7.}
  \label{fig:pva_6_7}
\end{figure}
